% Method: system overview first, then the measurement model, hardware and capture, panel experiments
% and the four applications. Every experiment has a continuous metric and a baseline (Table 1).
% The full method with equations, build detail and team notes sits in prose_backup/03_method.tex.
\section{Method}\label{sec:method}

\begin{itemize}
\item An Arduino draws each aperture pattern on a stripped LCD in front of the lens.
\item A laptop sends the patterns over serial and fires the camera over USB.
\item Table~\ref{tab:experiments} lists each experiment with its metric and baseline.
\end{itemize}

\subsection{Measurement model}\label{sec:model}

\begin{itemize}
\item A defocused point images as a scaled copy of the aperture. The blur grows with distance from the focal plane (Figure~\ref{fig:blur}).
\item The pattern flips across focus, so an asymmetric pattern tells near from far.
\item The panel loses light, leaks through black pixels and diffracts copies of each blur 31\,px away. A card mask has none of these faults.
\end{itemize}

\begin{figure}[!t]
\centering
\includegraphics[width=\columnwidth]{blur_model.png}
\caption{Predicted blur diameter }
\label{fig:blur}
\end{figure}

\subsection{Hardware and capture}\label{sec:hardware}

\begin{itemize}
\item The panel is a 1.8-inch colour LCD with its backlight removed~\cite{adafruit358}. A 3D-printed mount holds it, or a card mask, on a Canon EF 50\,mm f/1.8 lens on a Canon EOS 60D.
\item For each frame the laptop sends a pattern to the Arduino, waits for the liquid crystal to settle and fires the camera.
\item An open-disc frame at each end of a sequence shows any camera movement.
\item We laser-cut the same patterns from black card as the baseline (Figure~\ref{fig:patterns}).
\end{itemize}

\begin{figure}[!t]
\centering
\includegraphics[width=\columnwidth]{patterns.png}
\caption{Panel images at the $128 \times 160$ resolution the laptop sends.}
\label{fig:patterns}
\end{figure}

\begin{table}[!t]
\caption{Experiments, metrics and baselines. Every metric is continuous and carries a standard deviation over repeats. Card masks are laser-cut copies of the panel patterns. NCC is normalised cross-correlation, PSF point spread function, SPI serial peripheral interface, RMS root mean square, SSIM structural similarity index, SNR signal-to-noise ratio and MSE mean squared error.}
\label{tab:experiments}
\centering
\footnotesize
\renewcommand{\arraystretch}{1.15}
\begin{tabular}{@{}>{\raggedright}p{0.21\columnwidth}>{\raggedright}p{0.36\columnwidth}>{\raggedright\arraybackslash}p{0.33\columnwidth}@{}}
\toprule
Experiment & Metric & Baseline \\
\midrule
Pattern fidelity & NCC of PSF against design & Card mask of the same design \\
Transmission & Stops lost and leakage $\varepsilon$ & Card disc, 2.6-stop floor \\
Grid diffraction & Offset (px) and energy of the copies & Card disc, 31\,px prediction \\
Pair spectra & Fraction of spectral zeros the pair shares & Card pair, Lab~3 pair, ideal masks \\
Registration & Shift between bracketing open frames (px) & Card masks we swap by hand \\
Switching time & 10 to 90\,\% rise and fall time (ms) & SPI transfer and datasheet times \\
Depth & RMS depth error (\% of distance) & Card pair, single patterns \\
All-in-focus & SSIM against in-focus frames & Conventional frame at f/22 \\
Refocus & Edge width after refocus (px) & Open disc, f/5.6 frame, lens focus \\
Bokeh & NCC against the star by field angle & Card star, open disc \\
Flutter shutter & SNR gain (dB) from MSE against a stationary reference & Box and impulse exposures \\
\bottomrule
\end{tabular}
\end{table}

\subsection{Panel characterisation}\label{sec:panel}

\begin{itemize}
\item We image a point source at five distances from 700 to 1300\,mm through panel and card.
\item We measure light loss, leakage, grid copies and how well each blur matches its design.
\item We measure camera drift between frames and the time the panel takes to switch.
\end{itemize}

\subsection{Applications}\label{sec:apps}

\begin{itemize}
\item \textbf{Depth:} two complementary coded patterns~\cite{zhou2009pairs} give depth and a sharp image. We compare against a hand-swapped card pair and each pattern on its own.
\item \textbf{Refocus:} a $3 \times 3$ scan of the aperture lets us refocus after capture~\cite{ng2005lightfield}. We compare edge width against one frame through a smaller aperture.
\item \textbf{Custom bokeh:} a star-shaped aperture shapes out-of-focus highlights. One fixed shape suits a card mask, so we expect the card to give the closer match.
\item \textbf{Flutter shutter:} we switch the panel on and off within one exposure to make motion blur invertible~\cite{raskar2006coded}. We compare against an ordinary exposure of the same length in dim light.
\end{itemize}
